\documentclass[10pt, aspectratio=169]{beamer}

% --- TEMA Y ESTILO ---
\usetheme{Madrid}
\usecolortheme{whale}

% --- PAQUETES ---
\usepackage[utf8]{inputenc}
\usepackage[spanish, es-tabla]{babel}
\usepackage{amsmath, amssymb}
\usepackage{siunitx}
\usepackage{booktabs}
\usepackage{pgfpages}

% Descomentar la siguiente línea para imprimir un PDF con las notas del orador:
% \setbeameroption{show notes on second screen}

\sisetup{output-decimal-marker = {,}}

% --- INFORMACIÓN DEL CURSO ---
\title[Semiconductores y Diodos]{Física de Semiconductores y la Unión P-N}
\subtitle{De la Física del Silicio al Modelado de Diodos Reales}
\author[M.Sc. Bernardo Quiroga]{M.Sc. Ing. Bernardo Quiroga Turdera}
\institute[UCB Tarija]{Circuitos Electrónicos III (IMT-221) \\ UCB Tarija}
\date{11 de agosto de 2026}

\begin{document}

\begin{frame}
  \titlepage
\end{frame}

\begin{frame}{El Porqué en Mecatrónica y Biomédica}
  \begin{itemize}
      \item Los componentes pasivos ($R, L, C$) son \textbf{lineales y bilaterales}.
      \item Los semiconductores permiten controlar el flujo de corriente mediante campos eléctricos internos (\textbf{no-linealidad controlada}).
  \end{itemize}
  \vspace{0.3cm}
  \begin{alertblock}{Aplicación PFI Hito 1}
      ¿Cómo bloqueamos transitorios inductivos de motores o descargas estáticas en biosensores? Utilizando la \textbf{barrera de potencial de la unión P-N}.
  \end{alertblock}
  \note{Recuerden el Pitch de la sesión pasada: la electrónica analógica es la interfaz entre el mundo físico no lineal y el microcontrolador digital. Hoy entenderemos la física de la barrera que protegerá sus tarjetas.}
\end{frame}

\begin{frame}{Silicio Intrínsico y Dopaje (Extrínsico)}
  \begin{block}{Silicio Intrínsico}
      Red tetraédrica con enlaces covalentes. Concentración intrínseca $n_i \approx \SI{1.5e10}{\per\centi\meter\cubed}$ a \SI{300}{\kelvin} (27°C).
  \end{block}
  
  \begin{block}{Dopaje Tipo N (Donadores)}
      Inserción de átomos pentavalentes (Fósforo, Arsénico). Portadores mayoritarios: \textbf{Electrones libres} ($n \approx N_D$).
  \end{block}

  \begin{block}{Dopaje Tipo P (Aceptores)}
      Inserción de átomos trivalentes (Boro). Portadores mayoritarios: \textbf{Huecos} ($p \approx N_A$).
  \end{block}
  % Aquí el docente puede insertar \includegraphics con el esquema cristalino
\end{frame}

\begin{frame}{Corriente de Deriva vs. Difusión}
  \begin{columns}
      \column{0.5\textwidth}
      \begin{block}{1. Deriva (Drift)}
          Movimiento impulsado por un \textbf{campo eléctrico externo} ($\mathbf{E}$).
          \begin{equation*}
          J_{drift} = q (n \mu_n + p \mu_p) E
          \end{equation*}
      \end{block}
      
      \column{0.5\textwidth}
      \begin{block}{2. Difusión (Diffusion)}
          Movimiento impulsado por un \textbf{gradiente de concentración} (mayor a menor).
          \begin{equation*}
          J_{diff} = q D_n \frac{dn}{dx} - q D_p \frac{dp}{dx}
          \end{equation*}
      \end{block}
  \end{columns}
\end{frame}

\begin{frame}{Formación de la Unión P-N en Equilibrio}
  \textbf{Fenómeno Físico:}
  \begin{enumerate}
      \item Al unirse P y N, los electrones libres de N \textbf{difunden} hacia P y se recombinan con huecos.
      \item Esta recombinación deja \textbf{iones fijos} en la red cristalina: $N_D^+$ en N y $N_A^-$ en P.
      \item Se crea la \textbf{Zona de Carga Espacial (Deplexión)}, desprovista de portadores libres.
      \item \textbf{Resultado:} Aparece un Campo Eléctrico Interno ($\mathbf{E}_0$) que se opone a la difusión.
  \end{enumerate}
\end{frame}

\begin{frame}{Potencial de Contacto Integrado ($V_0$)}
  En equilibrio termodinámico, la corriente de difusión se iguala exactamente con la corriente de deriva ($J_{total} = 0$).
  
  \begin{block}{Ecuación del Potencial Integrado}
  \begin{equation*}
  V_0 = V_T \ln \left( \frac{N_A N_D}{n_i^2} \right)
  \end{equation*}
  \end{block}
  \vspace{0.2cm}
  Donde:
  \begin{itemize}
      \item $V_T = \frac{k_B T}{q} \approx \SI{25.85}{\milli\volt}$ a temperatura ambiente (\SI{300}{\kelvin}).
      \item Para el Silicio típico, $V_0 \approx \SI{0.6}{\volt} - \SI{0.8}{\volt}$.
  \end{itemize}
\end{frame}

\begin{frame}{La Unión P-N Bajo Polarización}
  \begin{columns}
      \column{0.5\textwidth}
      \textbf{Polarización Inversa ($V_D < 0$):}
      \begin{itemize}
          \item El potencial externo \textbf{refuerza} el campo interno $\mathbf{E}_0$.
          \item La zona de deplexión se ensancha.
          \item Fluye pequeña corriente de saturación inversa ($I_S$) térmica.
      \end{itemize}
      
      \column{0.5\textwidth}
      \textbf{Polarización Directa ($V_D > 0$):}
      \begin{itemize}
          \item El potencial externo \textbf{reduce} la barrera ($V_0 - V_D$).
          \item La zona de deplexión se estrecha.
          \item Mayoritarios difunden masivamente: la corriente crece exponencialmente.
      \end{itemize}
  \end{columns}
\end{frame}

\begin{frame}{Ecuación Característica de Shockley}
  \begin{block}{La Ecuación del Diodo Ideal}
  \begin{equation*}
  I_D = I_S \left( e^{\frac{V_D}{\eta V_T}} - 1 \right)
  \end{equation*}
  \end{block}
  
  \textbf{Parámetros Físicos:}
  \begin{itemize}
      \item $I_S$: Corriente de saturación inversa (sensible a la temperatura).
      \item $\eta$: Factor de idealidad ($1 \le \eta \le 2$).
      \item $V_T$: Voltaje térmico (\SI{25.85}{\milli\volt} a 27°C).
  \end{itemize}
\end{frame}

\begin{frame}{Resistencia Dinámica / Pequeña Señal ($r_d$)}
  La \textbf{Resistencia Dinámica ($r_d$)} es la pendiente inversa de la curva I-V evaluada en el punto de operación $Q(V_D, I_D)$:
  
  \begin{equation*}
  r_d = \left. \frac{d V_D}{d I_D} \right\vert{}_{Q} = \frac{\eta V_T}{I_D + I_S} \approx \frac{\eta V_T}{I_D}
  \end{equation*}
  
  \begin{exampleblock}{Ejemplo Rápido}
  Si $I_D = \SI{1}{\milli\ampere}$ y $\eta = 1$:
  \begin{equation*}
  r_d = \frac{\SI{25.85}{\milli\volt}}{\SI{1}{\milli\ampere}} \approx \SI{25.85}{\ohm}
  \end{equation*}
  \end{exampleblock}
\end{frame}

\begin{frame}{Taller Computacional en Python (Asíncrono)}
  \textbf{Objetivo:} Modelar el comportamiento térmico y dinámico del diodo 1N4148 de la BOM del curso (No asumir idealidad de 0.7V).
  \vspace{0.3cm}
  
  \textbf{Reto en Jupyter Notebook:}
  \begin{itemize}
      \item Graficar $I_D$ vs. $V_D$ variando la temperatura de 0°C a 80°C.
      \item Calcular y graficar la resistencia dinámica $r_d$ en función de la corriente de polarización en DC.
      \item Analizar impacto térmico sobre $I_S$ y $V_T$.
  \end{itemize}
\end{frame}

\begin{frame}{Conclusiones y Próximos Pasos}
  \textbf{Conclusiones:}
  \begin{enumerate}
      \item La unión P-N crea una barrera de potencial no lineal gobernada por la ecuación de Shockley.
      \item La temperatura altera drásticamente a $I_S$ y $V_T$.
      \item La resistencia dinámica $r_d$ disminuye a medida que aumenta la corriente de polarización en DC.
  \end{enumerate}
  \vspace{0.4cm}
  \begin{alertblock}{Próxima Clase (Jueves 13 de agosto)}
  Rectificación, filtros capacitivos, sujetadores y regulación Zener de precisión.
  \end{alertblock}
\end{frame}

\end{document}